% EGA V, English translation by P. Blass, 1987 (cote 24): archive cote 24, site page 49 (PDF page 50).
% typescript. Draft compared: none.
% AI-transcribed by Claude Opus 5.5 (claude-opus-5-5, adaptive thinking, effort high), Message Batches API, batch ega24b (msgbatch_01RGJJeHdR96U16ZWFuUjtDx), 2026-11-13; not yet checked by eye.
\scanpage{24/49}
\mspage{49}

[Marginal note illegible ask A.G]\hand{*}

(ii) $Y$sing $\neq$ $\underline{Y}$sup sing (set theoretically speaking let us say
(ii ter) There exists an $x \in X(K)$ and a tangent hyperplane H to $X_{K}$
at x which is not osculating at x by which we understand precisely
that x is not supersingular for the section of $O_{X_{K}}$ (1) that defines
H...)

These conditions imply that $\underline{Y}$sup sing $\neq \phi$ [Fr illegible ask
A.G.] dim $\underline{Y}$sup sing $\leq r - 2$ so that the image of Tsup sing in $\underline{P}$v has a
codimension $\geq 2$, and they imply also (iii) dim $T = r - 1$, i.e. T
is of codimension one in $\underline{P}$v.

\uline{Proof}

The equivalence of (i) and (i leir) is trivial its
equivalence with (ii) is a trivial consequence of 8.7 b), finally
the equivalence of (ii) and of (ii ter) is practically the
definition of $\underline{Y}$sup sing. Evidently (i ter) $\Rightarrow$ (i) it remains to
prove that (i) $\Rightarrow$ (i ter). We may evidently suppose that K is
algebraically closed and we are reduced to prove (taking into
account the hypothesis (i)) that there exists an open set $U \neq \phi$
such that $\xi \in U$ (K) implies that there exists exactly one point of
$\underline{Y}$sing (K) over $\xi$. This will follow from 8.7 c) which implies more
precisely.

\uline{Corollary 8.9}

Suppose that condition (i) of 8.8 is satisfied and let U be
the open subset of T of the points where T is smooth over K. Then

\margin{* Ask AG}
